\section{Conclusion}
\label{sec:conclusion}

% Compact NeurIPS-style closing. Evidence: canonical macros; Results/Discussion/Limitations;
% paper/analysis.md; paper/conclusion.md. No new experiments. No invented citations.

We asked whether, on CIFAR-10 depth~20 with Option~A and a frozen post-activation training
recipe matched to ResNet-20-A, replacing the residual block by a full pre-activation unit
reduces best test error.
All primary comparisons are single-seed runs under that shared protocol
(Sections~\ref{sec:setup}--\ref{sec:limitations}).

First, residual learning is supported in this setting:
ResNet-20-A attains $\ErrResNet\%$ best test error versus $\ErrPlain\%$ for Plain-20
($\ResidualAdvantagePP$\,pp advantage; Section~\ref{sec:results-residual}).
Second, full PreAct-ResNet-20-A reaches $\ErrPreAct\%$, i.e.\
$\DeltaPreActMinusResNet$\,pp higher error than the frozen ResNet-20-A control;
\textbf{the hypothesized improvement was not supported under the frozen training recipe and
single-seed setting}
(Section~\ref{sec:disc-negative}).
Third, Ablation~A ($\ErrAblationA\%$) improves on full PreAct by
$\DeltaAblationAVsPreAct$\,pp in magnitude, while differing from ResNet-20-A by only
$\DeltaAblationAMinusResNet$\,pp---a small single-run gap that we do not treat as a
validated improvement.
Ablation~B fails with numerical instability (\AblationBStatus) and is not used as an
accuracy baseline.

These findings are scoped to CIFAR-10, depth~20, Option~A, the present implementation, the
frozen baseline recipe, and single-seed evaluation; they do not extend to ImageNet, deeper
networks, other datasets, retuned PreAct schedules, or multi-seed statistical conclusions.
More broadly, the study illustrates that architectural block changes should be interpreted
jointly with the optimization recipe, and that a carefully scoped negative result can shrink
the hypothesis space---motivating controlled recipe-rescue and multi-seed validation rather
than a blanket rejection of pre-activation residual units.
